\chapter{Neural Network Foundations: Backpropagation, Activations, Initialisation and Regularisation}
\companion{3Blue1Brown: But what is a neural network? (Deep Learning, Chapter 1)}{\ytid{aircAruvnKk}}

Deep learning was ``the most drilled section'' for the InfoEdge hire: vanishing and exploding gradients, weight decay, how a network
learns (backprop), preventing overfitting, Xavier vs He initialisation. The OA asked about the universal approximation theorem, the
number of hidden layers in a ``deep'' network, classification losses, ResNets for vanishing gradients, and RBMs/DBNs. This chapter
covers each from scratch, with a backprop computation you can do on paper.

\section{From a neuron to a network}
\subsection{One neuron}
A neuron computes a weighted sum of its inputs plus a bias, then applies a non-linear \textbf{activation} function:
\[ z = \mathbf w^\top\mathbf x + b,\qquad a = g(z). \]
\begin{example}[: one neuron]
$\mathbf x = (1, 2)$, $\mathbf w = (0.5, -1)$, $b = 1$: $z = 0.5 - 2 + 1 = -0.5$. ReLU: $a = \max(0, -0.5) = 0$. Sigmoid: $a = \sigma(-0.5) = 0.378$.
\end{example}
Logistic regression is exactly one neuron with a sigmoid.

\subsection{Layers}
A \textbf{layer} is many neurons reading the same inputs: $\mathbf a^{(l)} = g(W^{(l)}\mathbf a^{(l-1)} + \mathbf b^{(l)})$. Stacking layers gives a
\textbf{multilayer perceptron (MLP)}. The first layer reads raw features; \textbf{hidden layers} build intermediate features; the output layer produces
the prediction. A dense layer from $n_{in}$ to $n_{out}$ has $(n_{in} + 1)n_{out}$ parameters.
\begin{example}[: counting parameters]
MLP 784--128--10 (MNIST): $(784 + 1)\cdot128 + (128 + 1)\cdot10 = 100{,}480 + 1{,}290 = 101{,}770$.
\end{example}

\subsection{Why non-linear activations are essential}
Without them, layer after layer of $W\mathbf x + \mathbf b$ collapses into one linear map: $W_2(W_1\mathbf x + \mathbf b_1) + \mathbf b_2 = W'\mathbf x + \mathbf b'$.
A 100-layer linear network is just linear regression. Non-linearities let the network bend; e.g. XOR needs a hidden layer with a non-linearity.

\subsection{``Deep'' and the universal approximation theorem}
A network is usually called \textbf{deep} when it has more than one hidden layer (often many). The \textbf{universal approximation theorem}: a network with
a single hidden layer and enough units, using a non-polynomial activation (sigmoid, ReLU), can approximate any continuous function on a bounded region
to any accuracy. It says such a network \emph{exists}; it does not say how many units are needed (possibly exponentially many), nor that gradient descent will
find the weights, nor that it will generalise. Depth helps because deep networks can represent many functions with exponentially fewer units by reusing
intermediate features (edges $\to$ shapes $\to$ objects).

\section{Activation functions}
\begin{center}\small
\begin{tabularx}{\linewidth}{l l l X}
\toprule
Name & Formula & Derivative & Notes\\
\midrule
Sigmoid & $\frac{1}{1 + e^{-z}}$ & $\sigma(1-\sigma)\le0.25$ & outputs in $(0,1)$; saturates; not zero-centred; use at binary outputs\\
Tanh & $\frac{e^z - e^{-z}}{e^z + e^{-z}}$ & $1 - \tanh^2\le1$ & zero-centred; still saturates; RNN/LSTM default\\
ReLU & $\max(0, z)$ & 0 or 1 & cheap, no saturation for $z>0$; ``dying ReLU'' when $z<0$ always\\
Leaky ReLU & $\max(\alpha z, z)$ & $\alpha$ or 1 & fixes dying ReLU; $\alpha = 0.01$\\
GELU / SiLU & $z\Phi(z)$ / $z\sigma(z)$ & smooth & Transformers (BERT, GPT use GELU)\\
Softmax & $\frac{e^{z_k}}{\sum e^{z_j}}$ & & multiclass output layer\\
\bottomrule
\end{tabularx}
\end{center}
Leaky ReLU at $-2$ with $\alpha = 0.01$: $-0.02$.

\section{Loss functions}
Regression: MSE, MAE, Huber. Binary classification: sigmoid output + binary cross-entropy. Multiclass: softmax + categorical cross-entropy. Multi-label:
one sigmoid per label + BCE. Margin: hinge. Ranking/similarity: contrastive, triplet. The output activation and loss are chosen together as a
likelihood (Chapter 12). Classification losses: cross-entropy/log loss, hinge, exponential, focal. MSE, MAE, Huber are regression losses.

\section{Training: forward pass, loss, backward pass}
\begin{enumerate}
\item \textbf{Forward}: compute activations layer by layer and the loss.
\item \textbf{Backward (backpropagation)}: compute the gradient of the loss with respect to every weight using the chain rule, from the output back to the input.
\item \textbf{Update}: $W\leftarrow W - \eta\,\partial L/\partial W$ (or Adam).
\end{enumerate}

\subsection{The chain rule, in words}
If $L$ depends on $a$, which depends on $z$, which depends on $w$, then a tiny change in $w$ changes $z$ by $\frac{\partial z}{\partial w}$ times as much, which
changes $a$ by $\frac{\partial a}{\partial z}$ times that, which changes $L$ by $\frac{\partial L}{\partial a}$ times that:
\[ \frac{\partial L}{\partial w} = \frac{\partial L}{\partial a}\cdot\frac{\partial a}{\partial z}\cdot\frac{\partial z}{\partial w}. \]
Backprop applies this systematically, \emph{reusing} the shared factors: the gradient flowing into layer $l$ is computed once from the gradient of layer
$l+1$, so the cost of the backward pass is about the same as the forward pass (not one pass per weight).

\subsection{Backprop for a dense layer}
Let $\boldsymbol\delta^{(l)} = \partial L/\partial\mathbf z^{(l)}$ (the ``error'' at layer $l$'s pre-activations). Then
\[ \frac{\partial L}{\partial W^{(l)}} = \boldsymbol\delta^{(l)}\mathbf a^{(l-1)\top},\quad \frac{\partial L}{\partial\mathbf b^{(l)}} = \boldsymbol\delta^{(l)},\quad
\boldsymbol\delta^{(l-1)} = \big(W^{(l)\top}\boldsymbol\delta^{(l)}\big)\odot g'(\mathbf z^{(l-1)}). \]
At the output with softmax/sigmoid + cross-entropy, $\boldsymbol\delta = \hat{\mathbf y} - \mathbf y$.

\begin{example}[: full backprop by hand]
Network: input $x = 1$; hidden $h = \sigma(w_1x)$ with $w_1 = 0.5$; output $\hat y = w_2h$ with $w_2 = 2$; loss $L = \frac12(\hat y - y)^2$ with $y = 1$.
\textbf{Forward}: $z_1 = 0.5$, $h = \sigma(0.5) = 0.622$, $\hat y = 1.245$, $L = \frac12(0.245)^2 = 0.030$.
\textbf{Backward}: $\partial L/\partial\hat y = 0.245$. $\partial L/\partial w_2 = 0.245\cdot h = 0.152$.
$\partial L/\partial h = 0.245\cdot w_2 = 0.490$. $\partial h/\partial z_1 = h(1-h) = 0.235$. $\partial L/\partial z_1 = 0.115$. $\partial L/\partial w_1 = 0.115\cdot x = 0.115$.
\textbf{Update} ($\eta = 0.1$): $w_2 = 2 - 0.0152 = 1.985$; $w_1 = 0.5 - 0.0115 = 0.4885$. Both moves reduce $\hat y$ toward 1.
\end{example}

\section{Vanishing and exploding gradients}
\subsection{Why they happen}
The gradient at an early layer is a \emph{product} of many factors, one per later layer: each factor involves a weight matrix and an activation derivative.
\begin{itemize}
\item If the factors are typically less than 1 in size, the product shrinks exponentially: \textbf{vanishing gradients}; early layers barely learn. Sigmoid's
  derivative is at most 0.25: through 8 sigmoid layers the factor is at most $0.25^8\approx1.5\times10^{-5}$.
\item If typically greater than 1, it grows exponentially: \textbf{exploding gradients}; updates become huge, loss jumps to NaN. A factor of 1.5 per step over 20
  RNN time steps: $1.5^{20}\approx3{,}325$.
\end{itemize}
RNNs are the extreme case (the same weight matrix multiplied at every time step, Chapter 20).

\subsection{Fixes}
\begin{center}\small
\begin{tabularx}{\linewidth}{X X}
\toprule
Fix & Why it works\\
\midrule
ReLU-family activations & derivative 1 on the active side: no shrinking from activations\\
Careful initialisation (Xavier, He) & keeps signal and gradient variance roughly constant across layers\\
Batch / Layer normalisation & keeps pre-activations in a well-scaled range, out of saturation\\
Residual (skip) connections & gradient flows through the identity path: $\partial(x + F(x))/\partial x = 1 + F'$\\
LSTM/GRU gates (RNNs) & additive cell-state path with forget gate near 1\\
Gradient clipping & caps exploding gradients (by norm)\\
Smaller learning rate, warmup & avoids large unstable steps\\
\bottomrule
\end{tabularx}
\end{center}

\section{Weight initialisation: Xavier and He}
\subsection{Why not zeros or a constant?}
If all weights in a layer start equal, every neuron computes the same output and receives the same gradient: they stay identical forever (\textbf{symmetry}). Random
initialisation breaks symmetry. But the \emph{scale} matters: too large and activations saturate or explode; too small and the signal dies after a few layers.

\subsection{The variance argument}
For $z = \sum_{i=1}^{n_{in}}w_ix_i$ with independent, zero-mean $w_i$ and $x_i$: $\Var(z) = n_{in}\Var(w)\Var(x)$. To keep $\Var(z)\approx\Var(x)$ we need
$\Var(w) = 1/n_{in}$.
\begin{itemize}
\item \textbf{Xavier/Glorot} (tanh, sigmoid): balance forward and backward passes: $\Var(w) = \frac{2}{n_{in} + n_{out}}$. Uniform version:
  $U[-\sqrt{6/(n_{in}+n_{out})}, +\sqrt{6/(n_{in}+n_{out})}]$.
\item \textbf{He/Kaiming} (ReLU): ReLU zeroes half the inputs, halving the variance, so double it: $\Var(w) = \frac{2}{n_{in}}$.
\end{itemize}
\begin{example}[: initialisation numbers]
He normal for $n_{in} = 512$: std $= \sqrt{2/512} = 0.0625$. Xavier uniform for $n_{in} = 100$, $n_{out} = 50$: limit $\sqrt{6/150} = 0.2$.
\end{example}

\section{Residual connections (ResNet)}
\textbf{Degradation problem}: before ResNet, simply adding more layers to a deep network made even the \emph{training} error worse (not overfitting: an optimisation
difficulty). A residual block outputs $\mathbf x + F(\mathbf x)$: the layers learn a correction to the identity. If the best thing is to do nothing, set $F = 0$, which is easy.
Gradients flow straight through the ``$+\mathbf x$'' path, so very deep networks (152 layers and more) train well. Transformers use the same idea around every sublayer.

\section{Normalisation layers}
\begin{itemize}
\item \textbf{Batch Normalisation}: for each feature, normalise across the mini-batch to mean 0, variance 1, then scale and shift with learnable $\gamma$, $\beta$:
  $\hat x = (x - \mu_B)/\sqrt{\sigma_B^2 + \epsilon}$, $y = \gamma\hat x + \beta$. At test time use running averages of $\mu$, $\sigma^2$. Allows higher learning
  rates, reduces sensitivity to initialisation, mild regulariser. Depends on batch size; awkward for RNNs and tiny batches.
  \begin{example}[: BatchNorm by hand]
  Batch $(1, 2, 3)$: $\mu = 2$, $\sigma^2 = 2/3$, $\sigma = 0.816$. For $x = 3$: $\hat x = 1.225$; with $\gamma = 2$, $\beta = 1$: $y = 3.449$. A BatchNorm layer on 64
  features has $2\times64 = 128$ trainable parameters.
  \end{example}
\item \textbf{Layer Normalisation}: normalise across features within each example: independent of batch size; standard in Transformers and RNNs.
\end{itemize}

\section{Regularisation in deep learning}
Covered conceptually in Chapter 1; specifics:
\begin{itemize}
\item \textbf{Dropout}: during training, zero each unit with probability $p$; \textbf{inverted dropout} scales the survivors by $1/(1-p)$ so the expected activation is
  unchanged and nothing needs to change at test time (dropout off). With $p = 0.2$, a kept activation of 5 becomes $6.25$. Interpretation: trains an ensemble of
  thinned networks sharing weights; prevents co-adaptation.
\item \textbf{Weight decay / L2}, \textbf{early stopping}, \textbf{data augmentation}, \textbf{label smoothing} (with $\epsilon = 0.1$ and 10 classes, the target for the
  true class becomes $0.9 + 0.01 = 0.91$ and $0.01$ for each other), \textbf{batch norm}, smaller models, transfer learning.
\end{itemize}

\section{Learning rate and getting unstuck}
See Chapter 5. In deep nets: warmup, cosine or step decay, reduce-on-plateau; momentum or Adam to traverse plateaus; restarts; check for dead ReLUs (fraction of units
always 0); always first verify the model can overfit one small batch.

\section{RBMs and Deep Belief Networks (historical)}
A \textbf{Restricted Boltzmann Machine} is a two-layer generative energy-based model (visible and hidden units, connections only between layers), trained with
contrastive divergence. A \textbf{Deep Belief Network} stacks RBMs, trained greedily layer by layer, then fine-tuned; it was how deep networks were first trained
successfully (2006) before ReLU, better initialisation and GPUs made direct training work. Today mostly of historical interest, but the OA asked about them.

\begin{practical}[: neural networks at InfoEdge]
The Tools slide lists ANNs, PyTorch and Keras/TensorFlow, and ``TensorFlow optimisation''. Typical uses: two-tower recommendation and matching networks; deep
click models combining embeddings of seeker, job and context; and fine-tuned Transformers. Practical lessons that interviewers like to hear: start with a
GBM baseline on tabular data; use embeddings for high-cardinality IDs; monitor gradient norms; use mixed precision and early stopping; serve with quantised or
distilled models for latency.
\end{practical}

\stuck{3Blue1Brown: Backpropagation, intuitively}{\ytid{Ilg3gGewQ5U}}
\section{Interview questions}

\begin{qa}{\pyq{INT: how does a network learn / backprop} Explain how a neural network learns, from the forward pass to the weight update.}
\oneline{Inputs flow forward through weighted sums and non-linearities to a prediction; a loss measures the error; backpropagation applies the chain rule from the
output back to every weight to get gradients, reusing intermediate results; and an optimiser nudges each weight against its gradient.}
\why
\textbf{Forward.} Each layer computes $\mathbf a^{(l)} = g(W^{(l)}\mathbf a^{(l-1)} + \mathbf b^{(l)})$; store $\mathbf z$ and $\mathbf a$ for the backward pass.
\textbf{Loss.} Cross-entropy or MSE.
\textbf{Backward.} Start with $\boldsymbol\delta^{(L)} = \partial L/\partial\mathbf z^{(L)}$ ($=\hat{\mathbf y} - \mathbf y$ for softmax + CE). For each layer going backward:
weight gradient $\boldsymbol\delta^{(l)}\mathbf a^{(l-1)\top}$, then pass the error back through the weights and the activation derivative.
\textbf{Update.} $W\leftarrow W - \eta\nabla_W L$ on a mini-batch; repeat for many batches and epochs.
Backprop is efficient because each layer's gradient reuses the next layer's: total cost about 2--3 forward passes, not one per weight.
\worked The hand computation above: $\partial L/\partial w_2 = 0.152$, $\partial L/\partial w_1 = 0.115$, updated weights 1.985 and 0.4885.
\pushes \emph{``What is autograd?''} Frameworks record the computation graph during the forward pass and apply the chain rule automatically (reverse-mode
automatic differentiation). \emph{``Why reverse mode?''} One scalar loss and millions of parameters: reverse mode gets all gradients in one backward pass.
\end{qa}

\begin{qa}{\pyq{INT: vanishing and exploding gradients} Why do vanishing and exploding gradients happen, and how do you fix each?}
\oneline{Because the gradient reaching early layers is a product of many per-layer factors (weights times activation derivatives); if those are mostly below 1 it shrinks
exponentially, above 1 it blows up; fixes are ReLU-type activations, proper initialisation, normalisation, residual connections, gated RNNs, gradient clipping and
smaller learning rates.}
\why
$\frac{\partial L}{\partial\mathbf a^{(1)}} = \prod_{l=2}^{L}\big(W^{(l)\top}\text{diag}(g'(\mathbf z^{(l)}))\big)\frac{\partial L}{\partial\mathbf a^{(L)}}$ (schematically). Sigmoid
derivatives $\le0.25$ make the product tiny. Large weights (or a recurrent matrix with largest singular value above 1) make it huge.
\textbf{Symptoms.} Vanishing: early layers' weights barely change, training stalls; gradient norms per layer fall by orders of magnitude toward the input.
Exploding: loss spikes or NaN; gradient norm spikes.
\textbf{Fixes.} Vanishing: ReLU/GELU, He/Xavier init, BatchNorm/LayerNorm, residual connections, LSTM/GRU, shorter paths (attention). Exploding: gradient clipping,
lower learning rate, proper init, normalisation, weight decay.
\worked 8 sigmoid layers: at most $0.25^8 = 1.5\times10^{-5}$. RNN with factor 1.5 over 20 steps: about 3{,}325.
\pitfall Saying ``ReLU completely solves vanishing gradients'': it solves the activation part, not the weight-product part, and introduces dead units.
\end{qa}

\begin{qa}{\pyq{OA and INT: Xavier vs He initialisation} What are Xavier and He initialisation, and why does He use 2 in the numerator?}
\oneline{Both set random weight variance by layer size to keep activation and gradient variances stable across layers: Xavier uses $2/(n_{in} + n_{out})$ for tanh/sigmoid,
He uses $2/n_{in}$ for ReLU, the 2 compensating for ReLU zeroing half of its inputs.}
\why $\Var(z) = n_{in}\Var(w)\Var(x)$. For a symmetric zero-mean input, ReLU passes half the mass, so $\E[\text{ReLU}(z)^2] = \frac12\Var(z)$; to keep the variance, double
$\Var(w)$: $2/n_{in}$. Xavier averages $n_{in}$ and $n_{out}$ to balance forward and backward.
\worked He normal $n_{in} = 512$: std $0.0625$. Xavier uniform $n_{in} = 100$, $n_{out} = 50$: limit $0.2$.
\pushes \emph{``What happens with zero init?''} Symmetry: all units in a layer stay identical. \emph{``Biases?''} Usually zero (fine, symmetry is broken by weights).
\end{qa}

\begin{qa}{\pyq{OA: ResNets and vanishing gradients} How do residual connections help train very deep networks?}
\oneline{Each block outputs $\mathbf x + F(\mathbf x)$, so the gradient has a direct identity path backward ($1 + \partial F/\partial\mathbf x$), which does not vanish, and each block
only needs to learn a small correction, which makes identity mappings easy and solves the degradation problem.}
\why Without skips, a 56-layer plain network had higher training error than a 20-layer one (optimisation failure). With skips, depth helps. Unrolled, a ResNet is like
an ensemble of many paths of different lengths.
\end{qa}

\begin{qa}{\pyq{OA: universal approximation theorem} What does the universal approximation theorem say and not say?}
\oneline{A one-hidden-layer network with enough units and a suitable non-linear activation can approximate any continuous function on a compact set arbitrarily well;
it does not say how many units are needed, that training will find the weights, or that the result will generalise.}
\why Depth gives exponential efficiency for compositional functions; this is why we use deep, not just wide, networks.
\end{qa}

\begin{qa}{\pyq{OA: number of hidden layers in a deep network} When is a network called ``deep''?}
\oneline{Conventionally when it has more than one hidden layer, i.e. at least two hidden layers between input and output; modern networks have tens to hundreds.}
\end{qa}

\begin{qa}{\pyq{OA: classification losses} Which of these are classification losses: cross-entropy, hinge, MSE, Huber, focal, exponential?}
\oneline{Cross-entropy, hinge, focal and exponential; MSE and Huber are regression losses.}
\why Focal loss $-(1-p)^\gamma\log p$ down-weights easy examples (object detection, imbalance). Exponential loss is AdaBoost's. Hinge for $y = +1$, score 0.3: $0.7$.
\end{qa}

\begin{qa}{\pyq{OA: RBMs and DBNs} What are RBMs and deep belief networks?}
\oneline{An RBM is a two-layer undirected generative model with visible and hidden units connected only across layers, trained by contrastive divergence; a DBN stacks RBMs
trained greedily layer by layer, historically used to pre-train deep networks before modern techniques made direct training possible.}
\end{qa}

\begin{qa}{\pyq{INT: dropout} How does dropout work at train and test time, and why does it regularise?}
\oneline{During training each unit is zeroed with probability $p$ and survivors are scaled by $1/(1-p)$ (inverted dropout); at test time dropout is off and no scaling is needed;
it regularises by preventing co-adaptation and approximately averaging an ensemble of thinned networks.}
\worked $p = 0.2$, kept activation 5: $5/0.8 = 6.25$.
\pushes \emph{``Dropout with BatchNorm?''} Their interaction can shift variance between train and test; place dropout after BN or use it mainly in the classifier head.
\emph{``Monte Carlo dropout?''} Keep dropout on at test time and average many passes to estimate uncertainty.
\end{qa}

\begin{qa}{\pyq{INT: weight decay; learning rate adjustment} How do weight decay and learning-rate schedules interact in deep learning training?}
\oneline{Weight decay shrinks weights each step by a factor proportional to the learning rate (for SGD), so schedules change effective regularisation; with Adam, use AdamW so
decay is decoupled from adaptive scaling; typical recipes pair warmup and cosine decay with weight decay 0.01--0.1.}
\end{qa}

\begin{qa}{Why do we need non-linear activations? What happens to a 10-layer network without them?}
\oneline{Without non-linearities, compositions of linear layers are a single linear map, so a 10-layer network can represent nothing a linear model cannot (it cannot learn XOR).}
\end{qa}

\begin{qa}{Compare sigmoid, tanh and ReLU. What is a dying ReLU?}
\oneline{Sigmoid and tanh saturate (gradients near 0 for large $|z|$), sigmoid is also not zero-centred; ReLU does not saturate for positive inputs and is cheap, but a unit whose
input is always negative outputs 0 with 0 gradient forever: a dead ReLU.}
\why Causes of dying: large learning rate or large negative bias. Fixes: Leaky ReLU, ELU, GELU, He init, lower learning rate.
\end{qa}

\begin{qa}{Explain Batch Normalisation. What happens at inference?}
\oneline{It normalises each feature over the mini-batch to zero mean and unit variance and then applies a learnable scale and shift; at inference it uses running averages of
the training mean and variance, so predictions do not depend on the batch.}
\worked Batch $(1,2,3)$, $\gamma = 2$, $\beta = 1$: output for 3 is $3.449$.
\pushes \emph{``Why does it help?''} Keeps activations in a good range, smooths the loss landscape, allows larger learning rates. \emph{``BN vs LayerNorm?''} LN normalises
across features per example; used in Transformers and RNNs; independent of batch size.
\end{qa}

\begin{qa}{\deep Why can't we initialise all weights to the same non-zero value?}
\oneline{Because every unit in a layer would compute the same output and receive the same gradient, so they would remain identical after every update and the layer would behave
like a single unit.}
\end{qa}

\begin{qa}{\deep Is a neural network with a sigmoid output and cross-entropy loss just logistic regression?}
\oneline{Its last layer is exactly a logistic regression, but on learned features from the hidden layers rather than raw inputs; with no hidden layers it is logistic regression.}
\why This is the cleanest way to explain what hidden layers do: learn a representation in which the classes become linearly separable.
\end{qa}

\begin{qa}{Count the parameters of an MLP 784--128--10, and of a BatchNorm layer on 64 features.}
\oneline{101{,}770 and 128 trainable (plus 128 running statistics that are not trained).}
\end{qa}

\begin{qa}{Gradient of binary cross-entropy with respect to the logit when $p = 0.8$, $y = 1$?}
\oneline{$p - y = -0.2$.}
\end{qa}

\begin{qa}{\deep Your training loss is flat from the first epoch. List what you check.}
\oneline{That the model can overfit a single small batch (if not, a bug), that labels match inputs, that the loss and output activation are consistent, that gradients are non-zero
and not NaN, that the learning rate is in a sensible range, that inputs are normalised, and that layers are not frozen.}
\why Common culprits: applying softmax before \code{CrossEntropyLoss} (which expects logits), forgetting \code{optimizer.zero\_grad()}, data loader returning the same
batch, a dead network from bad init.
\end{qa}

\begin{qa}{\deep Deep learning vs gradient-boosted trees for a tabular dataset with 5 lakh rows and 60 features. Which would you start with?}
\oneline{Gradient-boosted trees: on tabular data with heterogeneous features they usually match or beat neural networks with far less tuning; I would consider a neural network
if there are high-cardinality IDs or text/image inputs that benefit from learned embeddings, or to ensemble with the GBM.}
\end{qa}

\begin{qa}{What is label smoothing? Compute the targets for $\epsilon = 0.1$ with 10 classes.}
\oneline{Replace the one-hot target with $(1-\epsilon)$ on the true class plus $\epsilon/K$ spread over all classes: true class $0.91$, others $0.01$ each; it discourages
over-confident logits and improves calibration.}
\end{qa}

\begin{qa}{\deep A 3-layer linear chain $h_1 = w_1x$, $h_2 = w_2h_1$, $\hat y = w_3h_2$ with all weights 0.5 and $x = 1$. What is $\partial\hat y/\partial w_1$, and what happens with 50
such layers?}
\oneline{$\partial\hat y/\partial w_1 = w_3w_2x = 0.25$; with 50 layers the corresponding product is $0.5^{49}\approx1.8\times10^{-15}$: a vanishing gradient caused by weights alone.}
\end{qa}
